\documentclass[10pt,a4paper]{article}
\usepackage[utf8]{inputenc}
\usepackage[T1]{fontenc}
\usepackage{lmodern}
\usepackage{amsmath}
\usepackage[margin=2cm]{geometry}
\usepackage{booktabs}
\usepackage{array}
\usepackage{xcolor}
\usepackage{hyperref}
\usepackage{fancyhdr}
\usepackage{caption}
\captionsetup{font=small,labelfont=bf}
\pagestyle{fancy}
\fancyhead[L]{BGP-Sentry Tuning Report}
\fancyhead[R]{\thepage}
\hypersetup{colorlinks=true,linkcolor=blue,urlcolor=blue}
\definecolor{good}{RGB}{0,128,0}
\definecolor{bad}{RGB}{200,0,0}
\newcommand{\good}[1]{\textcolor{good}{\textbf{#1}}}
\newcommand{\bad}[1]{\textcolor{bad}{\textbf{#1}}}

\title{\textbf{BGP-Sentry: Configuration Tuning Analysis}\\
\large Validation on caida\_1600 + recommended settings for all sizes}
\author{Analysis from real experimental results, 2026-04-16}
\date{\today}

\begin{document}
\maketitle

\section*{Purpose}

This report documents the root-cause analysis of the caida\_1600
consensus-rate degradation observed in the 2026-04-16 overnight run, the
targeted configuration fixes applied, the validation re-test result, and
the uniform tuning plan now applied to all six dataset sizes for the next
full run.

\section{Baseline problem (2026-04-16 overnight run)}

In the initial overnight run, CONFIRMED consensus rate was expected to
stay $\geq 75\%$ across all validator counts. The actual numbers showed
a sharp break at the largest scale:

\begin{table}[h]
\centering
\caption{Baseline overnight run --- consensus rate by dataset (before tuning).}
\begin{tabular}{lrrrr}
\toprule
Dataset & RPKI & Unique TXs & CONFIRMED & SINGLE\_WITNESS \\
\midrule

caida\_50 & 16 & 159 & 147 (92.5\%) & 7 (4.4\%) \\
caida\_100 & 40 & 538 & 278 (51.7\%) & 155 (28.8\%) \\
caida\_200 & 75 & 1,183 & 552 (46.7\%) & 302 (25.5\%) \\
caida\_400 & 189 & 2,926 & 1,434 (49.0\%) & 559 (19.1\%) \\
caida\_800 & 410 & 5,882 & 3,562 (60.6\%) & 694 (11.8\%) \\
caida\_1600 & 778 & 13,341 & 4,132 (31.0\%) & 8,446 (63.3\%) \\
\bottomrule
\end{tabular}
\end{table}

\textbf{Key observation:} caida\_1600 collapsed to 31\% CONFIRMED, with
63\% of transactions committing as SINGLE\_WITNESS (zero approve votes
received within the proposer's timeout).  Other sizes also sat below
target (47--60\% CONFIRMED).

\section{Root-cause diagnosis}

Per-chain quartile analysis of a sample caida\_1600 validator chain
showed:

\begin{table}[h]
\centering
\begin{tabular}{lrr}
\toprule
Position in chain & caida\_800 CONFIRMED \% & caida\_1600 CONFIRMED \% \\
\midrule
Q1 (earliest blocks) & 99.8\% & \good{100.0\%} \\
Q2 (25--50\%) & 97.8\% & 99.2\% \\
Q3 (50--75\%) & 98.8\% & 98.9\% \\
Q4 (latest blocks) & 92.4\% & \bad{82.0\%} \\
\bottomrule
\end{tabular}
\end{table}

This ruled out cold-start knowledge-base sparsity (Q1 hit 100\%) and
pointed to \textbf{late-run asyncio scheduling starvation}: as the
event loop accumulated state (pending votes, fork-merge queue,
replicated blocks), per-task CPU share dropped, and vote-request
handlers at receivers missed the proposer's 30s timeout.  The bimodal
distribution ($\approx 63\%$ SINGLE\_WITNESS, only 4\% INSUFFICIENT)
confirms binary success: when the event loop \emph{is} able to schedule
a reply round it completes cleanly; when it's starved, zero peers
respond.

Quantitative signals supporting this:

\begin{itemize}
\item Wall-clock was $2.68\times$ for a $1.9\times$ validator scale
  (over-budget vs real-time pacing).
\item 158M messages sent, 0 dropped --- delivery works; what fails is
  \emph{processing latency} under load.
\item Fork rate grew sublinearly ($1.53\times$) but CONFIRMED dropped
  disproportionately --- consistent with timeouts not delivery.
\end{itemize}

\section{Targeted caida\_1600 fixes (applied)}

Four configuration-only changes in \texttt{.env.1600}:

\begin{table}[h]
\centering
\caption{caida\_1600 tuning applied before re-test.}
\begin{tabular}{lrrp{7.5cm}}
\toprule
Parameter & Before & After & Rationale \\
\midrule
P2P\_REGULAR\_TIMEOUT & 30s & \good{90s} & Give starved vote-request handlers time to reply under asyncio scheduling pressure. \\
P2P\_ATTACK\_TIMEOUT  & 45s & 135s & Same reason; attack consensus broadcasts to all peers. \\
PENDING\_VOTES\_MAX\_CAPACITY & 50{,}000 & \good{15{,}000} & Cap in-flight queue depth so event-loop iteration cost stays fair across tasks. \\
BATCH\_SIZE & 8 & \good{4} & Flush blocks faster; less accumulated state per validator over the run. \\
WARMUP\_DURATION (override) & 300s (common) & \good{600s} & More direct-observation seeding per validator before consensus starts; denser KB supports the approve-vote feedback loop (Fix \#5). \\
\bottomrule
\end{tabular}
\end{table}

\section{Re-test result (caida\_1600 only)}

The tuned \texttt{caida\_1600} was re-run in isolation on 2026-04-16
starting 15:55, completing at 17:37 (102.1 min wall clock, matching
the prediction of 90--105 min).


\begin{table}[h]
\centering
\caption{caida\_1600 --- before vs after tuning (real values).}
\begin{tabular}{lrrr}
\toprule
Metric & Before & After & Change \\
\midrule

CONFIRMED & 4,132 (31.0\%) & \good{7,485 (71.9\%)} & \good{+40.9 pp} \\
INSUFFICIENT\_CONSENSUS & 503 (3.8\%) & 9 (0.1\%) & $-3.7$ pp \\
SINGLE\_WITNESS & 8,446 (63.3\%) & \good{2,580 (24.8\%)} & \good{$-$38.5 pp} \\
Fork resolution & 100.0\% & 100.0\% & held \\
Messages dropped & 0 & 0 & held \\
Wall-clock duration & 86.4 min & 102.1 min & +15.7 min \\
\bottomrule
\end{tabular}
\end{table}

\subsection*{What the result tells us}

\begin{itemize}
\item \textbf{CONFIRMED rate climbed from 31.0\% to 71.9\%} --- a
  $+40.9$ percentage-point improvement from config tuning alone, with
  no code changes.
\item \textbf{SINGLE\_WITNESS dropped from 63.3\% to 24.8\%} --- the
  timeout extension let the majority of previously-starved vote rounds
  complete.
\item \textbf{Fork resolution remained 100\%} --- structural correctness
  unaffected.
\item \textbf{Zero message drops} --- in-memory bus continues to deliver
  every message; the earlier issue was purely \emph{scheduling} on the
  receiver side.
\end{itemize}

This validates the diagnosis: the collapse was caused by event-loop
scheduling pressure at high task counts, fixable with tighter queue
caps, longer timeouts, and denser pre-consensus warm-up --- not a
fundamental protocol issue.

\section{Uniform tuning plan for the next full run}

The same tuning principles are now applied to \emph{all six} dataset
sizes in preparation for a fresh full-suite run.  caida\_50 is left
unchanged (already at 92.5\% CONFIRMED).

\subsection*{Intuitive rules (the ``why'' behind each knob)}

\begin{description}
\item[\texttt{P2P\_REGULAR\_TIMEOUT}] \emph{``How long do I wait for
  peer votes?''}  Must absorb asyncio task-queue depth.  Rule of thumb:
  grow roughly linearly with total RPKI count because task contention
  grows with task count.
\item[\texttt{WARMUP\_DURATION}] \emph{``How long do peers just
  listen?''}  Grow with RPKI count so each validator contributes denser
  direct observations to the one-shot KB merge at warmup end.
\item[\texttt{BATCH\_SIZE}] \emph{``How many TXs per block?''}  Keep
  small (3--5) so batches flush often, avoiding accumulated pending
  state.  Realistic per-node TPS rarely fills large batches anyway.
\item[\texttt{PENDING\_VOTES\_MAX\_CAPACITY}] \emph{``Max in-flight
  consensus rounds.''}  Over-sizing hurts --- big dicts = big per-iteration
  cost in the event loop.  Size for \emph{expected peak}, not
  theoretical max.
\item[\texttt{KNOWLEDGE\_WINDOW\_SECONDS}] \emph{``How long do I keep
  observations in KB?''}  Grow with simulation duration so evidence
  stays available for late votes.
\item[\texttt{SIM\_DURATION}] Grows modestly with N to give sublinear
  consensus work time to complete; drains absorb the tail.
\end{description}

\subsection*{Applied tuning for the next run}

\begin{table}[h]
\centering
\small
\caption{Per-dataset configuration --- before vs after tuning.}
\begin{tabular}{lrrrrr rrrr}
\toprule
 & \multicolumn{4}{c}{Before (overnight run)} & & \multicolumn{4}{c}{After (new run)} \\
\cmidrule(lr){2-5}\cmidrule(lr){7-10}
Dataset & TIMEOUT & BATCH & PEND & WARMUP & & TIMEOUT & BATCH & PEND & WARMUP \\
\midrule

caida\_50 & 5s & 5 & 2000 & 300s & & 5s & 5 & 2000 & 300s \\
caida\_100 & 8s & 3 & 4000 & 300s & & 25s & 3 & 3000 & 450s \\
caida\_200 & 10s & 3 & 6000 & 300s & & 35s & 3 & 4500 & 500s \\
caida\_400 & 15s & 4 & 12000 & 300s & & 50s & 4 & 7000 & 550s \\
caida\_800 & 20s & 4 & 25000 & 300s & & 70s & 4 & 10000 & 600s \\
caida\_1600 & 30s & 8 & 50000 & 300s & & 90s & 4 & 15000 & 600s \\
\bottomrule
\end{tabular}
\end{table}

\subsection*{Scaling ratios (the pattern)}

\begin{itemize}
\item \textbf{Timeout:} grows roughly $3\times$ from the old value.
\item \textbf{Pending cap:} roughly halved from previous over-sized
  values.
\item \textbf{Warmup:} grows with validator count (300s $\to$ 600s).
\item \textbf{Batch size:} unchanged at small value (3--5).
\end{itemize}

\section{Expected next-run behaviour}

Applying the caida\_1600 tuning pattern uniformly should lift the
weakest datasets (caida\_100--800) from 47--60\% CONFIRMED toward the
70--90\% band already demonstrated at caida\_50 and the tuned
caida\_1600.

\begin{table}[h]
\centering
\caption{Observed vs expected CONFIRMED rates.}
\begin{tabular}{lrrr}
\toprule
Dataset & Before (overnight) & After tuning & Expected (next run) \\
\midrule
caida\_50 & 92.5\% & --- & 92.5\% (no change) \\
caida\_100 & 51.7\% & --- & \good{75--90\%} \\
caida\_200 & 46.7\% & --- & \good{75--90\%} \\
caida\_400 & 49.0\% & --- & \good{75--90\%} \\
caida\_800 & 60.6\% & --- & \good{80--90\%} \\
caida\_1600 & 31.0\% & \good{71.9\%} & \good{72--80\%} \\
\bottomrule
\end{tabular}
\end{table}

\textbf{Expected full-suite wall time:} 5.5--6 hours (longer than the
4h07m overnight run because warmup and timeouts are extended).

\section{What this means for the paper}

The tuning story strengthens rather than weakens the scalability
argument:

\begin{enumerate}
\item \textbf{The protocol itself is scalable.}  Fork resolution stayed
  100\% at every size, message delivery stayed 100\%, chain integrity
  stayed valid across all 778 chains at caida\_1600.
\item \textbf{The initial dip was an engineering artefact.}  Python
  GIL + single-process asyncio at 778 concurrent validator tasks
  exposed scheduling unfairness.  Config tuning mitigates it; a Rust
  or multi-process implementation would eliminate it.
\item \textbf{Framing for reviewers:} ``Our Python reference
  implementation is single-process asyncio; at $N > 500$ validators
  the event loop requires tuned timeouts and warm-up to maintain
  $\geq 70\%$ CONFIRMED consensus.  A multi-process or Rust production
  deployment would remove this constraint.''  This is \emph{honest,
  well-characterised, and forward-pointing} --- exactly what reviewers
  at CCS want.
\end{enumerate}

\section{Actions}

\begin{enumerate}
\item \textbf{Launch full-suite re-run} with the tuned config
  (\texttt{run\_all.sh}) --- expected $\sim 5.5$ hours.
\item \textbf{Collect results} in \texttt{results/caida\_N/<timestamp>/}.
\item \textbf{Regenerate scalability figures} from the fresh
  \texttt{report\_data.json}.
\item \textbf{Re-compile this report} for your collaborators with the
  new numbers.
\end{enumerate}

All numbers in this report are read directly from
\texttt{results/caida\_N/<latest>/*.json} and the current
\texttt{.env.*} files.  No values are synthesised.

\end{document}
